\section{Plugging the GPU EVM into Quasar Consensus}

\label{sec:cevm-consensus-integration}
%% ============================================================================

The previous section established that a single Metal/CUDA kernel can execute
a Cancun-era block end-to-end. This section answers a separate question:
\emph{where} in the Quasar consensus pipeline does that execution actually
get invoked, against \emph{which} interface, and at \emph{what} batch sizes
does the GPU win.

\subsection{What Quasar Decomposes Into}
\label{sec:cevm-quasar-recap}

The consensus repository (\texttt{github.com/luxfi/consensus} at SHA
\texttt{5bfc6ac}) is a pure-protocol library: it contains no execution
logic and no state model. It exposes two engine modes wired from the same
sub-protocols:

\begin{itemize}
  \item \textbf{Linear chain} (\texttt{protocol/nova}, driven by
    \texttt{engine/chain}) --- Photon $\to$ Wave $\to$ Focus $\to$ Ray $\to$
    Sink. Used by the C-Chain.
  \item \textbf{DAG} (\texttt{protocol/nebula}, driven by \texttt{engine/dag}
    and \texttt{protocol/field}) --- Photon $\to$ Wave (per frontier vertex)
    $\to$ Flare (cert/skip) $\to$ Horizon (safe-prefix scan) $\to$ Field
    (commit). This is the path described in \S\ref{sec:gpu} for the
    DAG-EVM.
\end{itemize}

The components, with file references in
\texttt{\textasciitilde/work/lux/consensus}:

\begin{itemize}
  \item \emph{DAG vertex production.} \texttt{protocol/field/field.go:30}
    defines \texttt{Proposer[V].Propose(ctx, parents) (V, error)}. The
    \emph{networking} layer pushes a payload into the proposer; the
    proposer returns a vertex ID and ships it. No bytecode runs here ---
    a vertex is just \texttt{(parents, author, round, payload bytes)}.
  \item \emph{Wave-based ordering.} \texttt{protocol/wave/wave.go} (the
    \texttt{Wave[T]} struct) drives threshold voting per frontier vertex.
    The FPC selector (\texttt{protocol/wave/fpc/}) randomises the
    threshold $\theta \in [\theta_{\min}, \theta_{\max}]$ from a PRF seed
    to defeat adaptive adversaries. Wave alone returns
    \texttt{(preferOK, confOK)} per round; it does not commit.
  \item \emph{FPC final commitment.} \texttt{protocol/field/field.go:107}
    \texttt{Driver[V].Tick} computes a \emph{safe prefix} via
    \texttt{computeSafePrefix} (line~119), then calls
    \texttt{Committer[V].Commit(ctx, ordered)}. The committer's
    \texttt{Commit} call is the post-FPC hook: every vertex passed in is
    finalised with all its causal ancestors finalised.
  \item \emph{Quasar certificate.} \texttt{protocol/quasar/types.go:40}
    \texttt{QuasarCert} carries the BLS aggregate, Corona share, and
    Z-Chain Groth16-rolled-up ML-DSA proof. The cert is produced
    \emph{after} \texttt{Commit} returns; it certifies the safe prefix,
    not individual transactions.
\end{itemize}

\subsection{Where Bytecode Execution Lives}
\label{sec:cevm-pipeline-position}

The crucial point: \emph{Quasar does not call into the EVM}. Quasar orders
opaque vertex IDs. The EVM call lives one layer above Quasar, in the VM
that hosts the chain.

\paragraph{Linear chain (Nova).}
\texttt{engine/chain/block/block.go:88--130} defines \texttt{ChainVM},
which the engine drives. The engine's call sequence
(\texttt{engine/chain/engine.go:947--1019}, function
\texttt{buildBlocksLocked}) is:

\begin{enumerate}
  \item \texttt{vm.BuildBlock(ctx)} returns a \texttt{block.Block} (an
    interface, not bytes; line~94 of \texttt{block.go}).
  \item \emph{Engine releases its mutex} (line~976) and calls
    \texttt{vmBlock.Verify(ctx)} (line~977). \textbf{This is the EVM
    execution hook.} \texttt{Verify} runs the bytecode and returns the
    post-state. If \texttt{Verify} fails, the block is dropped and never
    enters consensus.
  \item Engine adds the block to consensus
    (\texttt{t.consensus.AddBlock}, line~983) and self-votes.
  \item On finalisation: \texttt{vmBlock.Accept(ctx)}
    (line~1000 for K=1, line~877 for the multi-validator finaliser
    \texttt{tryFinalizeBlock}, and line~924 for the engine's
    \texttt{DrainAccepted} loop).
\end{enumerate}

So bytecode runs \emph{during vertex production}, before consensus voting.
\texttt{Accept} is a state-commitment hook --- it persists what
\texttt{Verify} already computed, it does not re-execute.

\paragraph{DAG mode (Nebula).}
The same split shifts one level: \texttt{Proposer.Propose} runs the
bytecode for the proposed vertex's transactions and returns a vertex ID
that already commits to the post-state root. \texttt{Field.Commit} then
delivers an \emph{ordered} list of vertex IDs whose post-state was
computed at proposal time. The committer's job is to \emph{apply} those
state diffs against the now-canonical history --- not to re-run the EVM.

This split is forced by the protocol shape: Quasar votes on vertex IDs,
and a vertex's ID must commit to its post-state for the cert to be
meaningful. We therefore execute speculatively at proposal, and apply
deterministically at commit. The Block-STM scheduler in \texttt{cevm}
(\texttt{lib/evm/gpu/parallel\_engine.cpp}) provides the speculative
re-execution-on-conflict that this requires.

\subsection{The Interface Quasar Calls}
\label{sec:cevm-interface}

The full \texttt{ChainVM} surface area Quasar's chain engine drives, from
\texttt{engine/chain/block/block.go:88--130} (verbatim):

\begin{lstlisting}[basicstyle=\ttfamily\footnotesize, language=go]
type ChainVM interface {
    Initialize(ctx context.Context, init Init) error
    BuildBlock(context.Context) (Block, error)
    ParseBlock(context.Context, []byte) (Block, error)
    GetBlock(context.Context, ids.ID) (Block, error)
    SetPreference(context.Context, ids.ID) error
    LastAccepted(context.Context) (ids.ID, error)
    GetBlockIDAtHeight(context.Context, uint64) (ids.ID, error)
    SetState(context.Context, uint32) error
    Shutdown(context.Context) error
    Version(context.Context) (string, error)
    Connected(context.Context, ids.NodeID, *version.Application) error
    Disconnected(context.Context, ids.NodeID) error
    HealthCheck(context.Context) (HealthCheckResult, error)
    NewHTTPHandler(context.Context) (http.Handler, error)
    WaitForEvent(context.Context) (Message, error)
}

type Block interface {
    ID() ids.ID
    Parent() ids.ID
    ParentID() ids.ID
    Height() uint64
    Timestamp() time.Time
    Status() uint8
    Verify(context.Context) error   // <-- bytecode here
    Accept(context.Context) error   // <-- commit state here
    Reject(context.Context) error
    Bytes() []byte
}
\end{lstlisting}

For the DAG path, \texttt{engine/dag/vertex/vertex.go:11--28} adds:

\begin{lstlisting}[basicstyle=\ttfamily\footnotesize, language=go]
type Vertex interface {
    ID() ids.ID
    Bytes() []byte
    Height() uint64
    Epoch() uint32
    Parents() []ids.ID
    Txs() []ids.ID
    Status() choices.Status
    Verify(context.Context) error
    Accept(context.Context) error
    Reject(context.Context) error
}
type DAGVM interface {
    ParseVertex(context.Context, []byte) (Vertex, error)
    BuildVertex(context.Context) (Vertex, error)
}
\end{lstlisting}

\textbf{Two hooks, both for the GPU EVM:} \texttt{Verify} for execution,
\texttt{Accept} for state commit.

\subsection{Concrete Integration Glue}
\label{sec:cevm-glue}

The \texttt{cevm} package
(\texttt{\textasciitilde/work/lux/cevm/cevm.go} at SHA \texttt{4d652b2})
already exports the C ABI:

\begin{lstlisting}[basicstyle=\ttfamily\footnotesize, language=go]
func ExecuteBlockV2(backend Backend, numThreads uint32,
    txs []Transaction) (*BlockResultV2, error)
\end{lstlisting}

The integration component is a Go struct that satisfies \texttt{block.Block}
and routes \texttt{Verify} into \texttt{cevm.ExecuteBlockV2}. The plumbing
is direct because the consensus library already keeps execution and
ordering decoupled --- the only state that flows from the cevm result
into the consensus layer is the 32-byte state root (carried in
\texttt{Block.Bytes()}) and the boolean \texttt{Verify} return.

\begin{lstlisting}[basicstyle=\ttfamily\footnotesize, language=go]
// cevmBlock satisfies engine/chain/block.Block and engine/dag/vertex.Vertex.
type cevmBlock struct {
    id        ids.ID
    parent    ids.ID
    height    uint64
    ts        time.Time
    txs       []cevm.Transaction
    backend   cevm.Backend
    threads   uint32
    // populated by Verify, read by Accept:
    result    *cevm.BlockResultV2
    stateRoot [32]byte
    raw       []byte
    state     uint8 // choices.Processing / Accepted / Rejected
    parent2   *cevmBlock // for state diff application at Accept
    db        kvWriter
}

func (b *cevmBlock) ID() ids.ID         { return b.id }
func (b *cevmBlock) Parent() ids.ID     { return b.parent }
func (b *cevmBlock) ParentID() ids.ID   { return b.parent }
func (b *cevmBlock) Height() uint64     { return b.height }
func (b *cevmBlock) Timestamp() time.Time { return b.ts }
func (b *cevmBlock) Status() uint8      { return b.state }
func (b *cevmBlock) Bytes() []byte      { return b.raw }

// Verify is the GPU EVM hook. Called by engine/chain/engine.go:977.
func (b *cevmBlock) Verify(ctx context.Context) error {
    if b.result != nil {
        return nil // memoised
    }
    res, err := cevm.ExecuteBlockV2(b.backend, b.threads, b.txs)
    if err != nil {
        return fmt.Errorf("cevm verify height=%d: %w", b.height, err)
    }
    // Cross-check vs the proposer's claimed root in raw block bytes.
    if !bytes.Equal(res.StateRoot[:], b.expectedRoot()) {
        return errStateRootMismatch
    }
    // Reject any tx with TxError; TxRevert is a valid in-protocol failure.
    for i, s := range res.Status {
        if s == cevm.TxError {
            return fmt.Errorf("tx %d returned TxError", i)
        }
    }
    b.result = res
    copy(b.stateRoot[:], res.StateRoot[:])
    return nil
}

// Accept is the state-commit hook. Called by engine/chain/engine.go:877,
// engine.go:924 (DrainAccepted), or engine.go:1000 (K=1 fast path).
func (b *cevmBlock) Accept(ctx context.Context) error {
    if b.result == nil {
        return errVerifyNotRun
    }
    // Persist the diff produced by Verify. The cevm bridge surfaces gas
    // and status; the underlying state diff is pulled from the host
    // bridge's account-state dictionary.
    if err := b.db.applyDiff(b.result, b.stateRoot); err != nil {
        return err
    }
    b.state = choicesAccepted
    return nil
}

func (b *cevmBlock) Reject(ctx context.Context) error {
    b.state = choicesRejected
    b.result = nil // free GPU result
    return nil
}
\end{lstlisting}

The VM-level \texttt{ChainVM.BuildBlock} constructs a \texttt{cevmBlock}
from the mempool, picks the backend (\texttt{cevm.AutoDetect()} at boot,
\texttt{cevm.Health()} verified), and returns it. The engine then drives
\texttt{Verify}, votes, and on quorum calls \texttt{Accept}. No code in
\texttt{consensus/} changes; the GPU EVM is a pure plug into
\texttt{block.ChainVM}.

\subsection{Where the Two Parallelisms Meet}
\label{sec:cevm-parallelism}

Quasar can finalise multiple DAG vertices concurrently --- specifically,
any antichain of the conflict-free DAG (\S\ref{sec:dag-evm}). The GPU EVM
can execute thousands of transactions concurrently within one launch.
These are \emph{orthogonal} dimensions:

\begin{table}[h]
\centering
\small
\begin{tabular}{lll}
\toprule
\textbf{Source} & \textbf{Granularity} & \textbf{Limit} \\
\midrule
Quasar DAG width    & vertices per round  & antichain size of DAG \\
GPU EVM batch size  & txs per dispatch    & GPU SIMD lane count \\
\bottomrule
\end{tabular}
\end{table}

The dimension that dominates is the one with the lower throughput floor.
On M1 Max the GPU launch cost is \SI{6}{\milli\second} cold,
\SI{4}{\milli\second} warm (\S\ref{sec:gpu:speedup}); the per-vertex
proposer cost on the consensus side is dominated by the wave round-trip
(\SI{250}{\milli\second} default \texttt{RoundTO},
\texttt{protocol/field/field.go:67}). \emph{Consensus dominates by
$\sim$40$\times$.} The GPU has plenty of slack to amortise launch cost
across a vertex's transactions, but it cannot make the wave round
shorter.

Pipelining helps when the GPU launch on vertex $V_n$ overlaps with the
wave round on $V_{n-1}$. In the linear chain mode, the engine signals
\texttt{t.signalPipeline()} after every accept (\texttt{engine.go:894}),
so a build/verify can begin while the previous block is still being
gossiped. Pipelining hurts when the proposer must wait for the GPU
launch to drain before issuing a fresh \texttt{Verify} --- the cevm
\texttt{ExecuteBlockV2} is currently synchronous (the C bridge blocks
on \texttt{[commandBuffer waitUntilCompleted]}). Async dispatch is
listed as an open item in \texttt{lib/evm/gpu/PERF.md} (priority~2).

\subsection{Saturation Math}
\label{sec:cevm-saturation}

The required regime, given the cevm dispatch crossover plot
(\texttt{lib/evm/gpu/PERF.md}, ``Crossover analysis''):

\begin{align*}
T_{\text{launch}}^{\text{cold}} &= \SI{6}{\milli\second} \\
T_{\text{launch}}^{\text{warm}} &= \SI{4}{\milli\second} \\
T_{\text{op}}^{\text{GPU}} &= \frac{1}{339\text{M ops/s}} \approx \SI{2.95}{\nano\second/op} \\
T_{\text{op}}^{\text{CPU}} &= \frac{1}{159\text{M ops/s}} \approx \SI{6.29}{\nano\second/op}
\end{align*}

Break-even (warm) for a generic workload of $W$ opcodes per dispatch:
\[
T_{\text{launch}}^{\text{warm}} + W \cdot T_{\text{op}}^{\text{GPU}}
< W \cdot T_{\text{op}}^{\text{CPU}}
\;\;\Rightarrow\;\;
W > \frac{T_{\text{launch}}^{\text{warm}}}{T_{\text{op}}^{\text{CPU}} - T_{\text{op}}^{\text{GPU}}}
\approx \frac{\SI{4}{\milli\second}}{\SI{3.34}{\nano\second}}
\approx 1.2 \times 10^{6} \text{ opcodes}.
\]

So a single dispatch needs $\sim$\num{1.2}M opcodes for the GPU to win
strictly on bytecode. The 100\,K txs/sec target translates as follows.
A vertex of 100 txs at 50 EVM ops/tx delivers \num{5000} ops per
dispatch --- well \emph{below} the break-even. A vertex of 100 txs at
\num{600} ops/tx (loop-heavy contracts) delivers \num{60000} ops ---
still below. The break-even at \num{1.2}M ops requires either:

\begin{enumerate}
  \item \emph{Fat workload per dispatch:} 100 txs $\times$ \num{12000}
    ops/tx (heavy contract --- DEX swap, AMM compound), or
  \item \emph{Bigger batch:} \num{12000} txs at \num{100} ops/tx, or
  \item \emph{Pipelined dispatch} that amortises launch across multiple
    blocks.
\end{enumerate}

The 100\,K txs/sec (\num{1000} vertices $\times$ \num{100} txs)
\emph{aggregate} is more than enough --- if 1000 vertices/sec is
emitted from a wave of 5--10 parallel proposers, each proposer dispatches
to a separate cevm host (\texttt{cevm.ExecuteBlockV2} is
goroutine-safe per \texttt{cevm.go:17}, with thread-local kernel
hosts). Aggregating across proposers gives \num{200}--\num{1000} txs per
dispatch on average, still below break-even at the \num{50}-op-per-tx
floor.

\textbf{The ecrecover and keccak paths are different.} Their break-even
sits much lower because the per-op CPU cost is high
(\SI{34}{\micro\second}/sig for ecrecover --- Table~\ref{tab:gpu-cancun}).
At \num{47000} signatures per dispatch the GPU runs in \SI{50}{\milli
\second} versus the CPU's \SI{1613}{\milli\second}. A 100-tx block with
two signatures per tx (200 sigs) is below crypto break-even too; a
\num{10000}-tx block (or \num{50} blocks of \num{200} sigs each
batched together) is firmly above it. \emph{In production, signature
verification is the path that justifies the GPU; opcode execution wins
only on fat workloads or via pipelining.}

\subsection{Limitations}
\label{sec:cevm-limitations}

\begin{itemize}
  \item \emph{No host callback at vertex boundary.}
    \texttt{cevm.ExecuteBlockV2} is a single C call that returns when the
    full batch finishes. The GPU kernel cannot ask Quasar mid-dispatch
    ``did this vertex finalise yet?'' because there is no host-callback
    primitive in the bridge. Consequence: the EVM cannot observe
    consensus state changes (e.g., a finalised cross-shard read) within
    a single \texttt{Verify}. The host bridge handles
    \texttt{CALL}/\texttt{CREATE} only via static analysis on the entry
    bytecode (\S\ref{sec:gpu:bridge}); cross-vertex callbacks would
    require either re-issuing a fresh dispatch or a CPU fallback path.
  \item \emph{Synchronous dispatch.}
    \texttt{[commandBuffer waitUntilCompleted]} blocks the calling
    goroutine. While the consensus engine has a per-block goroutine
    pool (\texttt{engine.go} \texttt{buildBlocksLocked}), a backed-up
    GPU draining a fat dispatch will pile up subsequent \texttt{Verify}
    calls. Async dispatch (priority~2 in PERF.md) is required for true
    pipelining across vertices.
  \item \emph{Account-state opcodes return defaults.}
    \texttt{BALANCE}, \texttt{EXTCODESIZE}, \texttt{EXTCODEHASH},
    \texttt{EXTCODECOPY}, \texttt{SELFBALANCE} read a host-supplied
    snapshot dictionary. For a vertex that touches an address outside
    the static-analysis closure, the kernel reads zero. The integration
    glue (\texttt{cevmBlock.Verify}) does not detect this and would
    accept a wrong-answer block; the parity-test corpus
    (\S\ref{sec:parity}) is the only catch. A \emph{stricter}
    integration would re-run on CPU when any opcode misses the
    dictionary --- not yet wired.
  \item \emph{Dynamic-target \texttt{CALL} forces CPU fallback.}
    Per \S\ref{sec:gpu:bridge}, a runtime-computed callee address
    routes the entire transaction to evmone on CPU. This is correct
    but breaks the GPU-batch invariant: a block of 100 txs in which
    one tx has a dynamic CALL splits into one CPU run plus a
    99-tx GPU run, doubling the total wall-clock for that block.
  \item \emph{Quasar does not vote on state roots, only vertex IDs.}
    A vertex's payload bytes commit to the post-state root, but
    Quasar itself never re-checks that root. If two honest validators
    execute the same vertex on different cevm backends (e.g., one
    Metal, one CPU-Parallel) and produce different roots due to a
    parity bug, only the four-way parity test corpus
    (\texttt{test/parity\_test.cpp}) catches it before deployment.
    Run-time disagreement would surface as a vote split, not a
    halt --- so divergent backends would silently lose finality on
    the affected branch. The fix path is exactly the parity corpus
    discipline: every backend must be byte-equivalent on every
    consensus-critical vector.
  \item \emph{No state-streaming for huge dispatches.} The cevm result
    is materialised entirely in host RAM (the V2 ABI returns a single
    \texttt{state\_root[32]} plus per-tx arrays). A 100\,K-tx dispatch
    that needed to stream state diffs to a database during execution
    cannot be expressed in the current ABI. Block-STM provides the
    speculation; persistence is bulk-applied at \texttt{Accept}. For
    bytecode that mutates millions of slots, this means the GPU result
    sits in RAM until the consensus quorum lands, which is acceptable
    at today's chain TPS but will need a streaming path before
    a single GPU dispatch routinely exceeds available host RAM.
\end{itemize}

\paragraph{Net.} The integration is small --- a \texttt{block.Block}
adapter, one C call per \texttt{Verify}, no consensus changes ---
because the consensus library is intentionally execution-free. The
hard work is on the GPU side (parity, host bridge, V2 kernel for
SIMD-cooperative dispatch). The right way to read this section is
that the seam is already correct; what remains is closing the
break-even gap so that GPU dispatch wins at vertex sizes the network
actually emits.

%% ============================================================================
